\documentclass[11pt]{article}
\usepackage[T1]{fontenc}
\usepackage[utf8]{inputenc}
\usepackage{amsmath,amssymb,amsthm,mathrsfs,geometry,hyperref}
\geometry{margin=1in}
\setlength{\emergencystretch}{3em}
\newtheorem{theorem}{Theorem}
\newtheorem{lemma}{Lemma}
\newcommand{\B}{\mathscr B}
\newcommand{\V}{\mathcal V}
\newcommand{\M}{\mathcal M}
\newcommand{\tr}{\operatorname{Tr}}
\title{Exact one-factor calculation of the weighted theta-source constants}
\author{}
\date{13 September 2026}
\begin{document}
\maketitle

This calculation finishes the finite moment reduction left explicit in
U0065/A2497 and U0066/A2559 of the retained transcript. It calculates the
weighted-source Grams controlling the circle length, and the derivative
Grams controlling its Fourier cutoff, on the original theta source.
It supplies no asymptotic bound on their growing-degree inverses.

\section{Original objects and domains}
Retain the source's actual finite packet polynomial
\(h(s)=\prod_{\rho\in Z}(s-\rho)^{m_\rho}\), with the full selected
orders, \(g=2\xi\), and its specified inverse Mellin representative
\(F_h\in\B\), where
\[
 \B=\{F\in C^\infty(\mathbb R_{>0}):
 \sup_{x>0}x^b|(-x\partial_x)^jF(x)|<\infty
 \text{ for every }b\in\mathbb Z,j\ge0\}.
\]
Write \(v_h(s)=\M F_h(s)=g(s)/h(s)\),
\(\M F(s)=\int_0^\infty F(x)x^s\,dx/x\), and
\(D=-x\partial_x\). For \(k\ge1\), set
\[
 c=k/2,\quad D^{(k)}=\sum_{j=1}^kD_j,\quad
 \V P=P(D^{(k)})F_h^{\otimes k},\quad S=\sum_{j=1}^ks_j.
 \tag{1}
\]
The polynomial source of degree at most \(N\) is \(\mathcal P_N\),
with the original basis \(1,S,\ldots,S^N\). Its arithmetic quotient
map is the unchanged remainder map
\(J_N:\mathcal P_N\to E=\mathbb C[S]/(\chi_{h,k})\).
Its full observation is
\[
 J^{(k)}\V=\eta_{h,k}J_N,\qquad
 \eta_{h,k}[P]=\upsilon_h^{\otimes k}P(A_h^{(k)})1,\qquad
 \upsilon_h=j_h(g/h).
 \tag{2}
\]
No coefficient of this unit is changed by the calculations below.

The original half-density coordinates are
\[
 r=\log x_k,\qquad z_i=\log x_i-\log x_k\ (i<k),
 \tag{3}
\]
and
\[
 (\mathcal U_kF)(r,\mathbf z)
 =e^{(kr+\sum_{i<k}z_i)/2}
 F(e^{r+z_1},\ldots,e^{r+z_{k-1}},e^r).
 \tag{4}
\]
The coordinate inverse in (3) has determinant one in logarithmic
coordinates. Consequently (4) is an isometry from
\(L^2((0,\infty)^k,d^kx)\) to
\(L^2(\mathbb R,dr;L^2(\mathbb R^{k-1},d\mathbf z))\), and
\(\partial_r\mathcal U_k=-\mathcal U_k(D^{(k)}-c)\).
These statements follow by substitution in the measure and by the chain
rule, respectively; for \(k=1\) the relative space is \(\mathbb C\).

\section{Weighted lines and exact finite moments}
For every real \(a\), define
\[
 w_a(t)=\frac{|v_h(1/2+a+it)|^2}{2\pi},\quad
 \mu_j(a)=\int_{\mathbb R}t^j w_a(t)\,dt,\quad
 \mathcal M_a(z)=\sum_{j\ge0}\mu_j(a)\frac{z^j}{j!}.
 \tag{5}
\]
The last expression is a formal series: each coefficient below uses
only finitely many of its entries. No analytic convergence assertion
about that series is needed.

All the integrals in (5) are finite. Indeed, for any real \(\sigma\)
and nonnegative integers \(b,j\),
\(y^b\partial_y^j(e^{\sigma y}F_h(e^y))\) decays exponentially at
both ends. To check this, expand the derivative into finitely many
terms \(e^{\sigma y}D^lF_h(e^y)\), use a positive integer exponent
larger than \(\sigma\) at \(+\infty\), and an integer exponent less
than \(\sigma\) at \(-\infty\) in the definition of \(\B\).
Any polynomial in \(y\) is absorbed by the resulting strict
exponential margins. Thus \(e^{\sigma y}F_h(e^y)\) is Schwartz.
Its Fourier transform with kernel \(e^{ity}\) is exactly
\(v_h(\sigma+it)\), which is therefore Schwartz on that line.

Let
\[
 m_{k,a}=w_0^{*(k-1)}*w_a,\qquad
 \mu^{(k,a)}_j=j![z^j]\mathcal M_0(z)^{k-1}\mathcal M_a(z),
 \tag{6}
\]
where the exponent zero in the case \(k=1\) means the identity for
convolution. Expansion of \((t_1+\cdots+t_k)^j\), with Fubini
justified by absolute polynomial moments, proves
\(\mu^{(k,a)}_j=\int t^jm_{k,a}(t)\,dt\).

\begin{theorem}[The actual weighted source Gram]
For \(P,Q\in\mathcal P_N\), the original weighted observation is
\[
 \begin{split}
 M_{k,N}(a)(P,Q)
 &:=\int_{(0,\infty)^k}x_k^{2a}
                 \overline{\V P(\mathbf x)}\V Q(\mathbf x)\,d^kx\\
 &=\int_{\mathbb R}\overline{P(c+a+iu)}Q(c+a+iu)
                 m_{k,a}(u)\,du.
 \end{split}
 \tag{7}
\]
In the unchanged basis \(1,S,\ldots,S^N\), its entries are exactly
\[
 [M_{k,N}(a)]_{b,d}
 =\sum_{r=0}^b\sum_{s=0}^d
 \binom br\binom ds(c+a)^{b+d-r-s}(-i)^ri^s
 (r+s)![z^{r+s}]\mathcal M_0^{k-1}\mathcal M_a.
 \tag{8}
\]
In particular these matrices require only \(\mu_j(0),\mu_j(a)\)
for \(0\le j\le2N\). Their masses are
\[
 [M_{k,N}(a)]_{0,0}=\mu_0(0)^{k-1}\mu_0(a).
 \tag{9}
\]
\end{theorem}
\begin{proof}
Put \(y_j=\log x_j\), use \(\sigma_j=1/2\) for \(j<k\) and
\(\sigma_k=1/2+a\), and apply ordinary Fourier Plancherel to
\(e^{\sum_j\sigma_jy_j}\V P(e^{\mathbf y})\). Integration by
parts has no boundary term by the estimates just proved. It sends each
\(D_j\) to multiplication by \(s_j=\sigma_j+it_j\). Hence its
Fourier transform is
\[
 P(c+a+i\textstyle\sum_jt_j)
 \prod_{j<k}v_h(1/2+it_j)\,v_h(1/2+a+it_k).
\]
The factor \((2\pi)^{-k}\) in Plancherel is exactly the product of
the \((2\pi)^{-1}\) factors in (5). Expanding the complex inner
product and integrating over the sum coordinate proves (7). The
binomial formula for each original monomial gives (8), including both
phases. Formula (9) is its constant entry. In particular the shift
\(c+a\) has not been replaced by \(c\), and the sum has not been
replaced by an average.
\end{proof}

Each matrix \(M_{k,N}(a)\) is positive definite. The entire nonzero
Mellin function \(v_h\) has isolated zeros on each vertical line, so
\(w_a>0\) almost everywhere. For \(k\ge2\), a convolution in (6)
is positive at every real sum: for almost every choice of the first
\(k-1\) variables all factors of its integrand are positive. For
\(k=1\), the almost-everywhere positivity suffices. A nonzero
polynomial has only finitely many zeros on the integration line in
(7), so its squared norm is strictly positive.

Let \(M=M_{k,N}(0)\). The constants required by A2497 and A2559 are
therefore the explicitly defined finite-matrix values
\[
 \kappa_{a,N}=\lambda_{\max}
       \bigl(M^{-1}(M_{k,N}(a)+M_{k,N}(-a))\bigr),\qquad a>0,
 \tag{10}
\]
where eigenvalues are understood in the positive inner product \(M\).
In particular there is a fully specified upper value avoiding an
additional extreme-eigenvalue calculation:
\[
 0<\kappa_{a,N}\le
 \mathfrak k_{a,N}:=
 \tr\bigl(M^{-1}(M_{k,N}(a)+M_{k,N}(-a))\bigr).
 \tag{11}
\]
To prove (11), conjugate by \(M^{1/2}\); the resulting matrix is
positive, and each eigenvalue is at most the sum of its eigenvalues.
The congruence is used to prove the inequality, not to replace the
source basis or its metric. With
\(\kappa_\pm=\lambda_{\max}(M^{-1}M_{k,N}(\pm a))\), the same
argument gives the individual upper values
\(\tr(M^{-1}M_{k,N}(\pm a))\).

\section{Derivative and Fourier-cutoff Grams, including the cross terms}
Write
\[
 a_h(t)=v_h(1/2+it)/\sqrt{2\pi},\qquad
 n(t)=|a_h'(t)|^2,\qquad
 b(t)=\overline{a_h'(t)}a_h(t).
 \tag{12}
\]
All three measures \(w_0(t)dt,n(t)dt,b(t)dt\) have finite absolute
polynomial moments by the same Schwartz argument. The last is a
complex measure; it is not replaced by a positive one. Define
\[
 \begin{array}{lll}
 m_k=w_0^{*k},&n_k=w_0^{*(k-1)}*n,&b_k=w_0^{*(k-1)}*b,\\
 \mathcal N(z)=\sum_{j\ge0}\int t^jn(t)dt\ z^j/j!,&&
 \mathcal C(z)=\sum_{j\ge0}\int t^jb(t)dt\ z^j/j!.
 \end{array}
 \tag{13}
\]
Let \(p\ge1\), \(R_P(S)=(S-c)^pP(S)\), and let primes on
\(R_P\) mean polynomial derivatives in the unchanged coordinate
\(S\). The derivative Gram in A2497/A2559 is
\[
 H_{p,N}(P,Q)=\int_{\mathbb R\times\mathbb R^{k-1}}
 (1+r^2)\overline{\partial_r^p\mathcal U_k\V P}\,
                    \partial_r^p\mathcal U_k\V Q\,dr\,d\mathbf z.
 \tag{14}
\]
\begin{theorem}[Complete derivative Gram]
With all polynomials in the following formula evaluated at \(c+iu\),
\[
 \begin{split}
 H_{p,N}(P,Q)=\int_{\mathbb R}\bigl(&
 \overline{R_P}R_Q(m_k+n_k)
 +\overline{R_P'}R_Q'm_k\\
 &+i\overline{R_P}R_Q'b_k
 -i\overline{R_P'}R_Q\overline{b_k}\bigr)(u)\,du.
 \end{split}
 \tag{15}
\]
Every entry is a finite expression in the moments encoded by
\(\mathcal M_0^k\), \(\mathcal M_0^{k-1}\mathcal N\), and
\(\mathcal M_0^{k-1}\mathcal C\), through degree \(2(N+p)\).
\end{theorem}
\begin{proof}
The chain rule in (4) gives
\(\partial_r^p\mathcal U_k\V P=(-1)^p\mathcal U_k\V R_P\).
The two occurrences of \((-1)^p\) cancel in (14). The term without
\(r^2\) is (7) at \(a=0\), with polynomial \(R_P\).
For the second term return to the independent logarithmic coordinates
\(y_1,\ldots,y_k\); its multiplier is exactly \(r=y_k\).
For the Fourier kernel \(e^{i\sum t_jy_j}\), multiplication by
\(y_k\) transforms to \(-i\partial_{t_k}\). The factor \(-i\)
cancels against its conjugate in the inner product. After retaining
the Plancherel constants in the product of the \(a_h\)'s, the
differentiated last factor is
\[
 a_h'(t_k)R_P(c+i\textstyle\sum_jt_j)
 +i a_h(t_k)R_P'(c+i\textstyle\sum_jt_j).
\]
Multiplication by its conjugate for \(P\) and by the displayed
expression for \(Q\) gives, respectively, the \(n_k\) term, the
derivative \(m_k\) term, and the cross terms
\(+i\overline{R_P}R_Q'b_k\) and
\(-i\overline{R_P'}R_Q\overline{b_k}\). Integrating the other
variables gives (15). Expansion in powers of the sum and Fubini prove
the stated coefficient products. The largest degree occurs in
\(\overline{R_P}R_Q\) and is at most \(2(N+p)\).
\end{proof}

\section{A finite choice of comparison parameters and its exact scope}
Set
\[
 \mathfrak h_{p,N}=\tr(M^{-1}H_{p,N})\ge0.
 \tag{16}
\]
Both sides in (16) are defined on the original coefficient space.
The same positive-matrix argument used for (11) gives
\(H_{p,N}\preceq\mathfrak h_{p,N}M\).
Here are the explicit parameters obtained from these completed
calculations. For prescribed positive error contributions
\(\eta_1,\eta_2\), choose
\[
 L\ge a^{-1}\log(1+\mathfrak k_{a,N}/\eta_1),
 \tag{17}
\]
and, for the twisted lattice used by all scalar holonomy phases, take an integer
\[
 J\ge\max\left\{2,\,
 1+\frac L{2\pi}
 \left(\frac{\mathfrak h_{p,N}}{(2p-1)\eta_2}\right)^{1/(2p-1)}
 \right\}.
 \tag{18}
\]
The ordinary lattice can omit the extra one; it has denominator
\(J\) instead of \(J-1\).

For completeness, the source estimates underlying this parameter
choice have direct proofs. Write \(\Psi=\mathcal U_k\V\) and
\(\mathcal C(s)=\int\Psi(r)^*\Psi(r+s)dr\). Weighted
Cauchy--Schwarz gives, for \(s>0\),
\[
 |v^*\mathcal C(s)v|
 \le e^{-as}\sqrt{v^*M_{k,N}(-a)v\ v^*M_{k,N}(a)v}.
\]
The negative shifts exchange the two weights. Unfolding one period
gives the Gram difference as the sum of the nonzero correlations
at shifts \(mL\); holonomy only multiplies each by a unit complex
number. Summing the geometric series and using \(2\sqrt{xy}\le x+y\)
proves its absolute quadratic-form bound
\((M_{k,N}(a)+M_{k,N}(-a))/(e^{aL}-1)\).
This is at most \(\eta_1M\) by (11) and (17).

Integration by parts \(p\) times in the retained Fourier transform,
then Cauchy--Schwarz with
\(\int(1+r^2)^{-1}dr=\pi\), gives
\(\|\widehat{\Psi v}(u)\|^2\le
\pi |u|^{-2p}v^*H_{p,N}v\).
For the ordinary lattice, summing \(|n|>J\) and bounding
\(\sum_{n>J}n^{-2p}\le J^{1-2p}/(2p-1)\) gives
\[
 0\preceq M(L)-M(L,J)
 \preceq\frac1{2p-1}\left(\frac L{2\pi J}\right)^{2p-1}H_{p,N}.
 \tag{19}
\]
The constants arise from \(\pi/L\) in each Fourier coefficient and
the two lattice directions; thus the factor in (19) is literal.
For phases \(0\le\theta<2\pi\),
\(|n+\theta/(2\pi)|\ge |n|-1\) for the omitted indices. Repeating
the integral bound replaces \(J\) by \(J-1\), and (18) makes this
at most \(\eta_2M\). With \(\eta_1+\eta_2<1\), the resulting
finite source Gram lies between
\((1-\eta_1-\eta_2)M\) and
\((1+\eta_1+\eta_2)M\).

For this quotient conclusion put \(q=\deg\chi_{h,k}\) and take
\(N\ge q-1\), so the unchanged remainder map is onto \(E\).
Taking the infimum over the unchanged affine sets
\(\{P:J_NP=v\}\) gives the same inequalities for the canonical
quotient Grams. Their sections differ by the original
\(\ker J_N=\chi_{h,k}\mathcal P_{N-q}\), with
\(\mathcal P_{-1}=0\), hence by the original theta
relations under (1)--(2). At a retained support label that correction
maps to its supported zero; it never becomes external absence.

These estimates compute the finite error parameters from the actual
one-factor shifted Mellin moments and their full derivative cross
terms. They do not bound \(\mathfrak k_{a,N}\) or
\(\mathfrak h_{p,N}\) uniformly as \(k,N\) grow. For the original
four endpoints, one may use \(N=2q_k+2\) in (17)--(18), exactly as
in A2497; all smaller polynomial spaces are restrictions of that
single source. The inverse \(M^{-1}\) and every original mass remain
visible, so this result cannot be mistaken for the missing
sub-threshold determinant estimate.

Finally, the weighted observations are norms of the already specified
representatives, not new maps identifying the arithmetic cohomology
with the Boolean stalk. In the primary Tau Base source the cohomology
is \(\mathsf A_\tau=R\Gamma(P,-)\), its right adjoint is
\(K_\tau(W)\simeq S_\eta W[1]\), and its actual top tensor quotient
is \(Q^{\otimes k}\). Equations (1)--(2) remain the same finite source
and full arithmetic observation into that construction. Multiplying
the source norm by \(x_k^{2a}\) has not been asserted to descend as
an endomorphism of \(Q^{\otimes k}\).

\end{document}
